% Source physical 717--719. Native dense two-column Chinese adaptation.
% DM prices are historical, not current prices.
\clearpage
\begingroup
\sffamily\fontsize{8}{10.5}\selectfont
\setlength{\parindent}{0pt}\setlength{\parskip}{3pt}
\setlength{\columnsep}{6mm}
\twocolumn[{\sffamily\Large 数学讲义\par}\vspace{3pt}\hrule\vspace{4pt}
{\sffamily 完整目录单张备索.}\par\vspace{8pt}]
\providecommand{\sgaCatalogEntry}[6]{\noindent\textbf{第 #1 卷:} #2#3. #4 页. #5. DM #6\par}
\sgaCatalogEntry{38}{R. Berger, R. Kiehl, E. Kunz, H.-J. Nastold, }{解析几何中的微分学}{IV, 134}{1967}{12,—}
\sgaCatalogEntry{39}{}{概率论讨论班 I}{II, 189}{1967}{14,—}
\sgaCatalogEntry{40}{J. Tits, }{单李群及其表示的表}{VI, 53}{1967}{6,80}
\sgaCatalogEntry{41}{A. Grothendieck, }{局部上同调}{VI, 106}{1967}{10,—}
\sgaCatalogEntry{42}{J. F. Berglund, K. H. Hofmann, }{紧半拓扑半群与弱殆周期函数}{VI, 160}{1967}{12,—}
\sgaCatalogEntry{43}{D. G. Quillen, }{同伦代数}{VI, 157}{1967}{14,—}
\sgaCatalogEntry{44}{K. Urbanik, }{预测理论讲义}{IV, 50}{1967}{5,80}
\sgaCatalogEntry{45}{A. Wilansky, }{泛函分析专题}{VI, 102}{1967}{9,60}
\sgaCatalogEntry{46}{P. E. Conner, }{周期映射讨论班}{IV, 116}{1967}{10,60}
\sgaCatalogEntry{47}{}{美国中西部范畴讨论班报告 I}{IV, 181}{1967}{14,80}
\sgaCatalogEntry{48}{G. de Rham, S. Maumary, M. A. Kervaire, }{挠率与简单同伦型}{IV, 101}{1967}{9,60}
\sgaCatalogEntry{49}{C. Faith, }{内射模与商环讲义}{XVI, 140}{1967}{12,80}
\sgaCatalogEntry{50}{L. Zalcman, }{解析容量与有理逼近}{VI, 155}{1968}{13,20}
\sgaCatalogEntry{51}{}{概率论讨论班 II}{IV, 199}{1968}{14,—}
\sgaCatalogEntry{52}{D. J. Simms, }{李群与量子力学}{IV, 90}{1968}{8,—}
\sgaCatalogEntry{53}{J. Cerf, }{关于三维球面的微分同胚 ($\Gamma_4=0$)}{XII, 133}{1968}{12,—}
\sgaCatalogEntry{54}{G. Shimura, }{自守函数与数论}{VI, 69}{1968}{8,—}
\sgaCatalogEntry{55}{D. Gromoll, W. Klingenberg, W. Meyer, }{整体 Riemann 几何}{VI, 287}{1968}{20,—}
\sgaCatalogEntry{56}{K. Floret, J. Wloka, }{局部凸空间理论引论}{VIII, 194}{1968}{16,—}
\sgaCatalogEntry{57}{F. Hirzebruch, K. H. Mayer, }{$O(n)$-流形、怪球与奇点}{IV, 132}{1968}{10,80}
\sgaCatalogEntry{58}{}{Riemann 曲面的 Kuramochi 边界}{IV, 102}{1968}{9,60}
\sgaCatalogEntry{59}{K. Jänich, }{可微 $G$-流形}{VI, 89}{1968}{8,—}
\sgaCatalogEntry{60}{G. S. Jones 编, }{微分方程与动力系统讨论班}{VI, 106}{1968}{9,60}
\sgaCatalogEntry{61}{}{美国中西部范畴讨论班报告 II}{IV, 91}{1968}{9,60}
\sgaCatalogEntry{62}{Harish-Chandra, }{半单李群上的自守形式}{X, 138}{1968}{14,—}
\sgaCatalogEntry{63}{F. Albrecht, }{控制理论专题}{IV, 65}{1968}{6,80}
\sgaCatalogEntry{64}{H. Berens, }{处理 Banach 空间上逼近过程的插值方法}{VI, 90}{1968}{8,—}
\sgaCatalogEntry{65}{D. Kölzow, }{测度的微分}{XII, 102}{1968}{8,—}
\sgaCatalogEntry{66}{D. Ferus, }{微分几何与微分拓扑中的全绝对曲率}{VI, 85}{1968}{8,—}
\sgaCatalogEntry{67}{F. Kamber, P. Tondeur, }{平坦流形}{IV, 53}{1968}{5,80}
\sgaCatalogEntry{68}{N. Boboc, P. Mustaţă, }{与二阶椭圆型线性微分算子相联系的调和空间}{VI, 95}{1968}{8,60}
\sgaCatalogEntry{69}{H. Bauer 编, }{位势论讨论班}{VI, 180}{1968}{14,80}
\sgaCatalogEntry{70}{M. H. Löb 编, }{逻辑暑期学校文集}{IV, 331}{1968}{20,—}
\sgaCatalogEntry{71}{}{Pierre Lelong 讨论班 (分析), 1967—1968 年}{VI, 190}{1968}{14,—}
\sgaCatalogEntry{72}{J. Barwise 编, }{无穷语言的句法与语义}{IV, 268}{1968}{18,—}
\sgaCatalogEntry{73}{J. F. Cerf, }{有限群作用讲义}{IV, 123}{1968}{10,—}
\sgaCatalogEntry{74}{A. Fröhlich, }{形式群}{IV, 140}{1968}{12,—}
\sgaCatalogEntry{75}{G. Lumer, }{函数代数与 Hardy 空间}{VI, 80}{1968}{8,—}
\sgaCatalogEntry{76}{R. G. Swan, }{代数 $K$ 理论}{IV, 262}{1968}{18,—}
\sgaCatalogEntry{77}{P.-A. Meyer, }{Markov 过程: Martin 边界}{IV, 123}{1968}{10,—}
\sgaCatalogEntry{78}{H. Herrlich, }{拓扑反射与余反射}{XVI, 166}{1968}{12,—}
\sgaCatalogEntry{79}{A. Grothendieck, }{加性余纤维化范畴与相对余切复形}{IV, 167}{1968}{12,—}
\sgaCatalogEntry{80}{B. Eckmann 编, }{三元组与范畴同调理论讨论班}{IV, 398}{1969}{20,—}
\sgaCatalogEntry{81}{J.-P. Eckmann, M. Guenin, }{统计力学中的代数方法}{VI, 131}{1969}{12,—}
\sgaCatalogEntry{82}{J. Wloka, }{基本空间与广义函数}{VIII, 131}{1969}{12,—}
\sgaCatalogEntry{83}{O. Zariski, }{代数曲面理论引论}{IV, 100}{1969}{8,—}
\sgaCatalogEntry{84}{H. Lüneburg, }{有限置换群的传递扩张}{IV, 119}{1969}{10,—}
\sgaCatalogEntry{85}{P. Cartier, D. Foata, }{交换与重排的组合问题}{IV, 88}{1969}{8,—}
\sgaCatalogEntry{86}{P. Hilton 编, }{范畴论、同调论及其应用 I}{VI, 216}{1969}{16,—}
\sgaCatalogEntry{87}{M. Tierney, }{稳定同伦论中的范畴构造}{IV, 65}{1969}{6,—}
\clearpage
\onecolumn
\endgroup
